\section{Grid structures and their automorphisms}\label{str:structures}\label{sec:structures}

We construct two families of finite structures with the same block sizes
and different total charges.  The structures encode edge states and the
locally compatible choices of those states at each vertex of a box.
Their automorphisms will serve two purposes: they will witness changes of
edge choices, and they will constrain the objects that a choiceless
computation can construct.  The construction and its central extension
are from~\cite[Sections~2--3]{cptCompanion2026}; we give the definitions
and proofs needed here.

All vector spaces in this section are over $\F$, and $1/2$ denotes $2\in\F$.
The fixed vocabulary consists of eight binary relations:
\[
 \tau=\{\Ed,\Cf,\EB,\VB,\Inc,\Zrel_0,\Zrel_1,\Zrel_2\}.
\]
The numerical subscripts denote elements of $\F$.

\subsection{The box and its edge states}

Fix an integer $L\ge2$.  Let $G_L=(V_L,E_L)$ be the box graph with
\[
 V_L=\{0,\ldots,L-1\}^3,
\]
whose edges join vertices differing by one in a single coordinate.
Orient every edge in the positive coordinate direction.  We often omit
$L$ from the notation, writing $V,E$ for $V_L,E_L$.  Define
$\epsilon_{ve}$ to be $1$ when $v$ is the head of $e$, $-1$ when it is the
tail, and $0$ otherwise.  The incidence map is
\[
 \partial:\F^E\longrightarrow\F^V,
 \qquad (\partial z)_v=\sum_{e\in E}\epsilon_{ve}z_e.
\]
Let $F=F_L$ be the set of elementary square faces.  A face in coordinate
directions $i<j$ is traversed from its lower corner in directions
$+i,+j,-i,-j$.  Its boundary vector $c^f\in\F^E$ has coefficient $1$ on
edges traversed in their positive orientation, $-1$ on edges traversed
against that orientation, and $0$ elsewhere.  At each vertex on the face,
one traversal edge enters and one leaves.  Consequently
\begin{equation}\label{str:face-boundary}
 \partial c^f=0\qquad(f\in F).
\end{equation}
Write $F(e)$ for the faces containing an edge $e$, and $F(v)$ for the faces
containing a vertex $v$.  If $d(v)$ is the degree of $v$, then
\begin{equation}\label{str:local-bounds}
 |F(e)|\le4,\qquad |F(v)|\le12,\qquad 3\le d(v)\le6.
\end{equation}
Indeed, an edge belongs to at most two squares in each of the two
transverse directions; for each pair of coordinate directions, a vertex
belongs to at most four squares.  Each coordinate supplies at least one
and at most two incident edges, including at the boundary when $L=2$.

On $\F^2$ use the alternating bilinear form
\[
 \omega((r,s),(r',s'))=rs'-sr'.
\]
An edge state on $e$ is an element of
\[
 P_e=(\F^2)^{F(e)}\times\F,
 \qquad (u,z)=((u_f)_{f\in F(e)},z).
\]
It is useful to regard these states as a group. We use the finite
Heisenberg multiplication over a field of odd characteristic
\cite[Section~2.1]{gurevichHadani2009}, with its alternating form
specified by the face-boundary signs. Put
\[
 B_e(u,u')=\frac12\sum_{f\in F(e)}c^f_e\omega(u_f,u'_f)
\]
and define
\begin{equation}\label{str:edge-product}
 (u,z)(u',z')=(u+u',z+z'+B_e(u,u')).
\end{equation}
Bilinearity yields
\[
 B_e(u,u')+B_e(u+u',u'')
 =B_e(u',u'')+B_e(u,u'+u''),
\]
which proves associativity.  The identity is $(0,0)$, and alternation gives
$(u,z)^{-1}=(-u,-z)$.  Thus the scalar coordinate of the group difference
$(u,z)^{-1}(u',z')$ is
\begin{equation}\label{str:edge-difference}
 z'-z-\frac12\sum_{f\in F(e)}c^f_e\omega(u_f,u'_f).
\end{equation}
This corrected difference will be recorded by the three relations
$\Zrel_\delta$.

\subsection{Atoms, relations, and charges}

For a charge vector $b=(b_v)_{v\in V}\in\F^V$, define the structure $A_b$
as follows.  Its edge block $Y_e$ contains one atom $(e,u,z)$ for every
$(u,z)\in P_e$.  Its configuration block $X_v$ contains one atom
\[
 t=\bigl(v,((u_e,z_e))_{e\ni v}\bigr)
\]
for every tuple of incident edge states satisfying
\begin{equation}\label{str:configuration}
 \begin{aligned}
 u_{e,f}&=u_{e',f}
 &&\text{for the two sides $e,e'$ of each $f\in F(v)$ at $v$},\\
 \sum_{e\ni v}\epsilon_{ve}z_e&=b_v.
 \end{aligned}
\end{equation}
The blocks are disjoint, and the two sorts are separately tagged.  Put
\[
 Y=\coprod_{e\in E}Y_e,\qquad
 X=\coprod_{v\in V}X_v,\qquad A_b=Y\mathbin{\dot\cup}X.
\]
These coordinates are labels for atoms; they are not set members of the
atoms.  In particular, the construction supplies no coordinate functions
or order as part of the input structure.

For $t\in X_v$ and $e\ni v$, let $y_e(t)\in Y_e$ denote the edge atom
specified by the $e$-coordinate of $t$.  Interpret the relations by
\[
 \begin{aligned}
 \Ed&=\{(y,y):y\in Y\},&
 \Cf&=\{(t,t):t\in X\},\\
 \EB&=\bigcup_{e\in E}Y_e^2,&
 \VB&=\bigcup_{v\in V}X_v^2,\\
 \Inc&=\{(t,y_e(t)):v\in V,\ t\in X_v,\ e\ni v\}.&&
 \end{aligned}
\]
For $\delta\in\F$, set
\begin{equation}\label{str:z-relation}
 \Zrel_\delta((e,u,z),(e,u',z'))
 \quad\Longleftrightarrow\quad
 z'-z-\frac12\sum_{f\in F(e)}c^f_e\omega(u_f,u'_f)=\delta.
\end{equation}
Each relation is false on all pairs not specified by these rules.
In particular, $\Inc$ is directed from configurations to edge atoms, and
$\Zrel_\delta$ only relates atoms in one edge block.  Exactly one of the
three $\Zrel_\delta$ holds for each ordered pair within an edge block.

\begin{lemma}\label{str:size}
For every $L\ge2$ and every $b\in\F^{V_L}$, all edge and configuration
blocks of $A_b$ are nonempty, and
\begin{equation}\label{str:size-formula}
 |A_b|=N_L=
 \sum_{e\in E_L}3^{2|F(e)|+1}
 +\sum_{v\in V_L}3^{2|F(v)|+d(v)-1}.
\end{equation}
Thus $N_L$ is independent of $b$, and
\[
 |Y_e|\le3^9,\qquad |X_v|\le3^{29},\qquad
 L^3\le N_L\le(3^{10}+3^{29})L^3.
\]
\end{lemma}
\begin{proof}
The edge-block count is immediate from the definition of $P_e$.
At a vertex $v$, each face $f\in F(v)$ contributes one freely chosen
vector in $\F^2$, shared by its two sides at $v$.  Coordinates for different
faces are distinct, so these choices contribute $3^{2|F(v)|}$ possibilities.
The scalar constraint in Equation~\eqref{str:configuration} is one equation
in $d(v)$ variables, with every coefficient nonzero.  Choose the scalars
on all but one incident edge arbitrarily; the remaining scalar is then
uniquely determined.  Hence it contributes $3^{d(v)-1}$ possibilities
for every value of $b_v$.  This proves Equation~\eqref{str:size-formula}
and nonemptiness.  The block bounds follow from
Equation~\eqref{str:local-bounds}.  Finally, there are $L^3$ vertices and
$3L^2(L-1)\le3L^3$ edges, which gives the asserted bounds on $N_L$.
\end{proof}

The quotient sets of $\EB$-blocks and $\VB$-blocks are therefore precisely
$\{Y_e:e\in E\}$ and $\{X_v:v\in V\}$.  Moreover, there is an $\Inc$-pair
from $X_v$ to $Y_e$ if and only if $e\ni v$: every configuration chooses
one state on each incident edge, and every $X_v$ is nonempty.  Thus the
incidence graph between these blocks recovers the underlying box.

The two families used below have charges
\begin{equation}\label{str:charges}
 b^0=0,\qquad b^1=\mathbf1_{v_*},\qquad v_*=(0,0,0).
\end{equation}
By Lemma~\ref{str:size}, $A_{b^0}$ and $A_{b^1}$ have the same size $N_L$.
Their total charges in $\F$ are respectively $0$ and $1$.  The next
construction gives a common family of automorphisms for both structures,
and indeed for every charge vector.

\subsection{A central extension acting on the full structure}

Define the vector spaces
\[
 C_F=(\F^2)^F,\qquad K=\ker\partial\subseteq\F^E.
\]
An element of $K$ is a circulation: its signed sum at every vertex is
zero.  For $a,a'\in C_F$, put
\[
 B(a,a')=\frac12\sum_{f\in F}\omega(a_f,a'_f)c^f.
\]
Equation~\eqref{str:face-boundary} implies $B(a,a')\in K$.
On the set $H=C_F\times K$, define multiplication by
\begin{equation}\label{str:group-law}
 (a,k)(a',k')=(a+a',k+k'+B(a,a')).
\end{equation}

\begin{lemma}\label{str:central-extension}
The multiplication in Equation~\eqref{str:group-law} makes $H$ a finite
group.  Its identity is $(0,0)$ and $(a,k)^{-1}=(-a,-k)$.
The subgroup $\{0\}\times K$, identified with $K$, is central, and
projection onto $C_F$ is a surjective homomorphism with kernel $K$.
With the convention $[h,h']=hh'h^{-1}(h')^{-1}$, one has
\begin{equation}\label{str:commutator}
 [(a,k),(a',k')]
 =\left(0,\sum_{f\in F}\omega(a_f,a'_f)c^f\right).
\end{equation}
\end{lemma}
\begin{proof}
The map $B$ is bilinear and alternating.  Its bilinearity gives
\[
 B(a,a')+B(a+a',a'')=B(a',a'')+B(a,a'+a''),
\]
so both associations of a product of three elements in $H$ agree.
The identity and inverse follow from $B(0,a)=B(a,0)=B(a,-a)=0$.
The assertions about the central subgroup and projection now follow
directly from the product formula.  Finally, multiplying the four factors
of the commutator gives $(0,B(a,a')-B(a',a))=(0,2B(a,a'))$.
Since $2(1/2)=1$ in $\F$, this is Equation~\eqref{str:commutator}.
\end{proof}

For $(a,k)\in H$, define a map $h_{a,k}$ on each edge block by
\begin{equation}\label{str:edge-action}
 h_{a,k}(e,u,z)=
 \left(e,\ (u_f+a_f)_{f\in F(e)},\
 z+k_e+\frac12\sum_{f\in F(e)}c^f_e\omega(a_f,u_f)\right).
\end{equation}
On a configuration atom, apply this map to each of its incident edge
states.  The following proposition verifies that the resulting tuple is
again a configuration in the same block and that all input relations are
preserved.

\begin{proposition}\label{str:action}
For every $L\ge2$ and $b\in\F^{V_L}$, the maps $h_{a,k}$ define a faithful
action of $H$ by automorphisms of the entire $\tau$-structure $A_b$.
Every $h_{a,k}$ preserves each edge and configuration block setwise.
For $k\in K$, the central action $h_{0,k}$ fixes all face vectors and adds
$k_e$ to the scalar coordinate on edge $e$.
\end{proposition}
\begin{proof}
For each edge $e$, the restriction map
\[
 \rho_e:H\longrightarrow P_e,
 \qquad \rho_e(a,k)=((a_f)_{f\in F(e)},k_e)
\]
is a homomorphism: the $e$-coordinate of $B(a,a')$ is exactly the
correction in the product on $P_e$.  Equation~\eqref{str:edge-action} is
left multiplication by $\rho_e(a,k)$.  It therefore defines a group
action on the edge states, with inverse $h_{-a,-k}$.

Fix a configuration at $v$.  For each $f\in F(v)$, write $u_f$ for its
common vector on the two incident sides of $f$.  Both copies receive the
same addition $a_f$, so the face equalities remain true.  The change in
its divergence is
\[
 (\partial k)_v+
 \frac12\sum_{f\in F(v)}
 \left(\sum_{e\ni v}\epsilon_{ve}c^f_e\right)\omega(a_f,u_f)=0.
\]
Here the first term vanishes because $k\in K$, and every inner sum
vanishes by Equation~\eqref{str:face-boundary}.  The image is consequently
a configuration with charge $b_v$.  The inverse map gives a bijection
of each $X_v$, and componentwise composition gives the same group law
as on edge atoms.

Preservation of sorts and blocks proves preservation of
$\Ed,\Cf,\EB,\VB$.  The coordinatewise definition on configurations
proves preservation of $\Inc$.  For the remaining relations, let
$g,g'\in P_e$ be two edge states and $p=\rho_e(a,k)$.  The group identity
\[
 (pg)^{-1}(pg')=g^{-1}g'
\]
shows that their corrected scalar difference in
Equation~\eqref{str:edge-difference} is unchanged.  Thus every
$\Zrel_\delta$ is preserved.  These statements also apply to the inverse,
so all eight relations are preserved in both directions, including
nonrelations.  This proves that the action is by automorphisms of $A_b$.

If $h_{a,k}$ fixes every atom, apply it to the state $(0,0)$ on each edge.
Equation~\eqref{str:edge-action} then gives $k_e=0$ for every $e$ and
$a_f=0$ whenever $f\in F(e)$.  Every face has an edge, so $a=0$ and $k=0$;
the action is faithful.  The final assertion follows by setting $a=0$
in Equation~\eqref{str:edge-action}.
\end{proof}

We identify $H$ and $K$ with these permutation groups.  Their actions
extend to $\HF(A_b)$ by the recursive rule
$h\cdot x=\{h\cdot y:y\in x\}$ when $x$ is a set, and by the specified
permutation when $x$ is an atom.  Every object has finite rank, so this
rule is well defined at every rank; it preserves membership and fixes
all pure hereditarily finite sets.

Two face-vector translations can have a nontrivial commutator in the
central circulation subgroup $K$. This link between the quotient and the
central subgroup will be used in Section~\ref{sec:support}: a small orbit
under the full group $H$ will constrain which circulations in $K$ can
move an object. For the choice construction, the two kinds of coordinates
will instead be corrected successively, first by a face translation and
then by a circulation.
